\documentclass[10pt]{article}
\usepackage[a4paper,margin=0.86in]{geometry}
\renewcommand{\rmdefault}{cmr}
\usepackage{amsmath,amssymb,amsthm,mathtools,amscd,booktabs,enumitem,xcolor,hyperref}
\hypersetup{colorlinks=true,linkcolor=blue!55!black,citecolor=blue!55!black,urlcolor=blue!55!black}
\setlength{\emergencystretch}{3em}
\setlength{\parindent}{0pt}
\setlength{\parskip}{0.52em}
\setlist{leftmargin=1.6em,itemsep=0.22em,topsep=0.22em}
\sloppy

\newtheorem{theorem}{Theorem}[section]
\newtheorem{proposition}[theorem]{Proposition}
\newtheorem{corollary}[theorem]{Corollary}
\newtheorem{lemma}[theorem]{Lemma}
\theoremstyle{definition}
\newtheorem{definition}[theorem]{Definition}
\newtheorem{convention}[theorem]{Convention}
\newtheorem{remark}[theorem]{Remark}

\newcommand{\Ql}{\mathbf Q_\ell}
\newcommand{\Qlb}{\overline{\mathbf Q}_\ell}
\newcommand{\Dbc}{\mathcal D^b_c}
\newcommand{\Dbcnorm}{\mathcal D^b_{c,\mathrm{norm}}}
\newcommand{\Mod}{\mathsf{Mod}}
\newcommand{\Comp}{\operatorname{Comp}}
\newcommand{\Kar}{\operatorname{Kar}}
\newcommand{\Hom}{\operatorname{Hom}}
\newcommand{\pred}[1]{\par\smallskip\begingroup\footnotesize\raggedright
\emph{Immediate predecessors:} #1.\par\endgroup}

\begin{document}
\begin{center}
{\large\bfseries Base Change for Compact Support\par}
\medskip
{\normalsize Finite torsion, finiteness, and passage to l-adic coefficients\par}
\end{center}

\begin{abstract}
For a separated finite-type morphism, compact-support base change is formed
from open-extension exchange and proper base change after choosing a
compactification.  Refinement compatibility makes the result independent
of that choice.  The finite-torsion construction is followed by
constructibility and cohomological-dimension theorems, then by integral
adic realization, rationalization, and extension to algebraically closed
l-adic coefficients.  All derived categories are ordinary
Verdier-localized triangulated categories.
\end{abstract}

\tableofcontents

\section{Finite-torsion categories and four elementary functors}

\begin{convention}[Geometric and coefficient hypotheses]
Let \(S\) be noetherian, let \(\ell\) be invertible on \(S\), and fix
\(\Lambda_n=\mathbf Z/\ell^n\mathbf Z\).  All schemes are noetherian and of
finite type over \(S\).  A morphism to which \(Rf_!\) is applied is separated
and of finite type.  The topology \(\tau\) is the small etale or small
pro-etale topology; finite constructible objects on the two sites are
identified by the finite-torsion comparison equivalence.
\end{convention}

\begin{definition}[Derived and constructible categories]
Let \(\Mod_\tau(T,\Lambda_n)\) be the Grothieck abelian category of sheaves
of \(\Lambda_n\)-modules on \(T_\tau\).  Its ordinary derived category is
\[
 D(T,\Lambda_n)=
 \operatorname{Cpx}(\Mod_\tau(T,\Lambda_n))
 [\{\text{quasi-isomorphisms}\}^{-1}].
\]
\(D^+(T,\Lambda_n)\) is the full triangulated subcategory of complexes
whose cohomology vanishes in all sufficiently negative degrees.
A sheaf is constructible if a finite constructible stratification of \(T\)
makes it locally constant with finite stalks on each stratum.  The full
triangulated subcategory
\[
 \Dbc(T,\Lambda_n)\subset D(T,\Lambda_n)
\]
consists of complexes with constructible cohomology sheaves, all zero
outside one finite interval; thus \(\Dbc(T,\Lambda_n)\subset
D^+(T,\Lambda_n)\).
\end{definition}

\begin{definition}[Derived extension of finite coefficients]
Let \(A\to B\) be a homomorphism of finite rings and let \(B_T\) be the
associated constant ring sheaf on \(T_\tau\).  Write
\(D(T,A)=D(\Mod_\tau(T,A))\), and similarly for \(B\).  For
\(K\in D(T,A)\), choose a K-flat replacement \(P_A(K)\to K\) and define
\[
 K\otimes_A^LB:=P_A(K)\otimes_AB_T
 \quad\in D(T,B).
\]
This is independent of the K-flat replacement up to a unique natural
isomorphism.  If \(B\) has finite Tor dimension over \(A\), it preserves
\(D^+\) and bounded constructible complexes.  For
\(A\to B\to C\), the derived tensor associator
\[
 (K\otimes_A^LB)\otimes_B^LC
 \xrightarrow{\sim}K\otimes_A^LC
\]
is fixed with its identity and threefold-associativity equations.
\end{definition}

\begin{definition}[Inverse image]
For \(h:T'\to T\), inverse image is exact and therefore gives
\[
 h^*:D(T,\Lambda_n)\longrightarrow D(T',\Lambda_n).
\]
It preserves \(D^+\) and \(\Dbc\).  The canonical maps
\((kh)^*\simeq h^*k^*\) and \(\operatorname{id}^*\simeq\operatorname{id}\)
are fixed with their associativity and unit equations.
\end{definition}

\begin{definition}[Derived direct image]
For \(p:T\to Y\), the sheaf direct image
\[
 p_*:\Mod_\tau(T,\Lambda_n)\to\Mod_\tau(Y,\Lambda_n)
\]
is left exact.  A functorial K-injective resolution \(K\to I(K)\) defines
\[
 Rp_*K=p_*I(K),\qquad
 Rp_*:D(T,\Lambda_n)\longrightarrow D(Y,\Lambda_n).
\]
It restricts to \(Rp_*:D^+(T,\Lambda_n)\to D^+(Y,\Lambda_n)\).
For composable \(p,r\), fix the natural composition isomorphism
\[
 \mu_{r,p}:Rr_*Rp_*\xrightarrow{\sim}R(rp)_*
\]
with its identity and associativity equations.  No constructible landing
assertion is included in this definition.
\end{definition}

\begin{definition}[Open extension by zero]
Let \(j:U\hookrightarrow T\) be a quasi-compact open immersion.  For a
sheaf \(F\) on \(U_\tau\), let \(j_!F\) be the sheafification of
\[
 V\longmapsto
 \{s\in F(V\times_TU):
   \operatorname{supp}(s)\text{ is closed in }V\}.
\]
This is left adjoint to \(j^*\).  For a geometric point \(\bar t\to T\),
\[
 (j_!F)_{\bar t}\simeq
 \begin{cases}
 F_{\bar t},&\bar t\in U,\\
 0,&\bar t\notin U.
 \end{cases}
\]
\end{definition}

\begin{theorem}[Exactness and constructibility of open extension]
The sheaf functor \(j_!\) is exact and carries finite constructible sheaves
to finite constructible sheaves.  Termwise application preserves
quasi-isomorphisms and gives exact triangulated functors
\[
 j_!:D(U,\Lambda_n)\to D(T,\Lambda_n),\qquad
 j_!:D^+(U,\Lambda_n)\to D^+(T,\Lambda_n),\qquad
 j_!:\Dbc(U,\Lambda_n)\to\Dbc(T,\Lambda_n).
\]
For composable open immersions, \((kj)_!\simeq k_!j_!\), coherently with
identities and threefold composition.
\end{theorem}
\pred{the support-closed definition, its stalk formula, and conservativity of geometric stalks}

\section{The two exchange morphisms}

\begin{definition}[Open-extension exchange]
For a Cartesian square
\[
\begin{CD}
 U'@>{j'}>>T'\\
 @V{a}VV @VV{b}V\\
 U@>{j}>>T
\end{CD}
\]
with \(j,j'\) quasi-compact open immersions, pullback of support-closed
sections defines
\[
 \beta_{j,b}:b^*j_!\longrightarrow j'_!a^*.
\]
It is first defined on sheaves and then applied termwise to complexes.
\end{definition}

\begin{theorem}[Open-extension exchange isomorphism]
The map \(\beta_{j,b}\) is a natural isomorphism on \(D(U,\Lambda_n)\).
It preserves bounded constructible objects and satisfies:
\[
 \beta_{j,\operatorname{id}}=\operatorname{id},
\]
the exchange for a composite pullback is the composite of the two
exchanges, and the exchange for a composite of open immersions agrees with
the two successive open exchanges.
\end{theorem}
\pred{exact open extension, exact inverse image, and pullback of closed supports}

\begin{definition}[Proper base-change morphism]
For a Cartesian square
\[
\begin{CD}
 T'@>{a}>>T\\
 @V{p'}VV @VV{p}V\\
 Y'@>{b}>>Y
\end{CD}
\]
with \(p\) proper, define
\[
 \operatorname{BC}_{p,b}:b^*Rp_*\longrightarrow Rp'_*a^*
\]
as the adjunction mate obtained from the unit of \(a^*\dashv Ra_*\), the
specified Cartesian identification \(pa=bp'\), the composition maps for
derived direct image, and the counit of \(b^*\dashv Rb_*\).
\end{definition}

\begin{theorem}[Finite-torsion proper base change]
For \(K\in D^+(T,\Lambda_n)\),
\[
 \operatorname{BC}_{p,b}(K):
 b^*Rp_*K\xrightarrow{\sim}Rp'_*a^*K
\]
is an isomorphism in \(D(Y',\Lambda_n)\).  It is natural in \(K\), equals
the identity for an identity square, and satisfies horizontal and vertical
pasting for Cartesian squares.  For a geometric point \(\bar y\to Y\),
\[
 (Rp_*K)_{\bar y}\simeq
 R\Gamma(T_{\bar y},K|_{T_{\bar y}}).
\]
\end{theorem}
\pred{the proper base-change morphism, the strict-henselian stalk comparison, and invariance of finite proper-fiber cohomology}

\begin{theorem}[Open--proper interchange]
For a Cartesian square
\[
\begin{CD}
 T'@>{j'}>>\overline T'\\
 @V{q'}VV @VV{q}V\\
 T@>{j}>>\overline T
\end{CD}
\]
in which \(j,j'\) are quasi-compact open immersions and \(q,q'\) are
proper, the mate of open-extension exchange and proper base change is a
natural isomorphism
\[
 \operatorname{Ex}_{j,q}:j_!Rq'_*\xrightarrow{\sim}Rq_*j'_!
\]
on \(D^+(T',\Lambda_n)\).  These maps satisfy identity, horizontal-pasting,
and vertical-pasting equations.
\end{theorem}
\pred{open-extension exchange, finite-torsion proper base change, and the adjunctions for open inverse image and proper direct image}

\section{Compactifications and the ambient functor
\texorpdfstring{\(Rf_!\)}{Rf!}}

\begin{definition}[Compactification and dense compactification]
A compactification of \(f:X\to Y\) is a factorization
\[
 c=(X\xrightarrow{j_c}\overline X_c\xrightarrow{p_c}Y)
\]
with \(j_c\) a quasi-compact open immersion, \(p_c\) proper, and
\(p_cj_c=f\).  It is dense if
\(\mathcal O_{\overline X_c}\to(j_c)_*\mathcal O_X\) is injective.
Write \(\Comp_Y(f)\) for all compactifications and
\(\Comp_Y^{\mathrm{dens}}(f)\) for the full subcategory of dense ones.
\end{definition}

\begin{definition}[Refinement and Cartesian refinement]
A refinement \(q:c_1\to c_2\) is a proper \(Y\)-map
\(q:\overline X_{c_1}\to\overline X_{c_2}\) satisfying
\[
 qj_{c_1}=j_{c_2},\qquad p_{c_2}q=p_{c_1}.
\]
It is Cartesian over the common open if the square
\[
\begin{CD}
 X@>{j_{c_1}}>>\overline X_{c_1}\\
 @V{\operatorname{id}}VV @VV{q}V\\
 X@>{j_{c_2}}>>\overline X_{c_2}
\end{CD}
\]
is Cartesian.  The source of a refinement dominates its target.
\end{definition}

\begin{theorem}[Dense replacement]
Every compactification \(c\) admits a canonical dense refinement
\[
 \delta_c:c^{\mathrm{dens}}\longrightarrow c
\]
obtained by taking the scheme-theoretic closure of \(X\) in
\(\overline X_c\).  A refinement whose source is dense is Cartesian over
the common open.
\end{theorem}
\pred{scheme-theoretic closure and the open-and-closed section criterion for a proper refinement}

\begin{theorem}[Common dense refinements]
Dense compactifications exist.  For arbitrary
\(c_1,c_2\in\Comp_Y(f)\), the scheme-theoretic closure of the image of
\[
 X\longrightarrow\overline X_{c_1}\times_Y\overline X_{c_2}
\]
is a common dense refinement of \(c_1,c_2\).  Parallel refinement maps are
equal after a further refinement, and in fact are already equal when their
source is schematically dense and their target is separated.  Hence the
full subcategory \(\Comp_Y^{\mathrm{dens}}(f)\) is cofiltered and is
coinitial in \(\Comp_Y(f)\) for arrows from dominating to dominated
compactifications.
\end{theorem}
\pred{Nagata compactification, scheme-theoretic closure, and stability of proper morphisms under products}

\begin{definition}[The compactification candidate]
For any \(c\in\Comp_Y(f)\), define the ambient functor
\[
 Rf_{!,c}:=Rp_{c,*}j_{c,!}:
 D^+(X,\Lambda_n)\longrightarrow D^+(Y,\Lambda_n).
\]
Only bounded-below derived functors enter this definition; constructible
landing will be stated only after the finiteness theorems.
\end{definition}

\begin{definition}[Refinement transformation]
For a Cartesian refinement \(q:c_1\to c_2\), open-proper exchange gives
\(\xi_q:j_{c_2,!}\xrightarrow{\sim}Rq_*j_{c_1,!}\).
Let \(\gamma_q=\xi_q^{-1}\), and define
\[
 \theta_q:
 Rp_{c_1,*}j_{c_1,!}
 \xrightarrow{\mu^{-1}}
 Rp_{c_2,*}Rq_*j_{c_1,!}
 \xrightarrow{Rp_{c_2,*}\gamma_q}
 Rp_{c_2,*}j_{c_2,!}.
\]
\end{definition}

\begin{theorem}[Cartesian-refinement coherence]
Each \(\theta_q\) for a Cartesian refinement is a natural isomorphism of functors
\(D^+(X,\Lambda_n)\to D^+(Y,\Lambda_n)\), and
\[
 \theta_{\operatorname{id}}=\operatorname{id},\qquad
 \theta_{rq}=\theta_r\theta_q.
\]
\end{theorem}
\pred{open-proper exchange and coherence of derived direct image}

\begin{definition}[Comparison of arbitrary compactifications]
For \(c_1,c_2\in\Comp_Y(f)\), choose a common dense refinement
\(r_i:d\to c_i\) and define
\[
 \kappa^d_{c_1,c_2}:=
 \theta_{r_2}\theta_{r_1}^{-1}:
 Rf_{!,c_1}\xrightarrow{\sim}Rf_{!,c_2}.
\]
\end{definition}

\begin{theorem}[Independence of the common dense refinement]
The map \(\kappa^d_{c_1,c_2}\) is independent of \(d\) and its refinement
arrows.  Denote it by \(\kappa_{c_1,c_2}\).  It satisfies
\[
 \kappa_{c,c}=\operatorname{id},\qquad
 \kappa_{c_2,c_3}\kappa_{c_1,c_2}=\kappa_{c_1,c_3}.
\]
\end{theorem}
\pred{common dense refinements and Cartesian-refinement coherence}

\begin{definition}[Chosen ambient exceptional direct image]
Choose \(c_f\in\Comp_Y^{\mathrm{dens}}(f)\) and set
\[
 Rf_!:=Rf_{!,c_f}:
 D^+(X,\Lambda_n)\longrightarrow D^+(Y,\Lambda_n).
\]
Changing \(c_f\) uses the natural isomorphism \(\kappa\) from refinement
coherence.
\end{definition}

\section{Finite compact-support exchange through one compactification}

\begin{definition}[Raw base-changed compactification]
Fix a Cartesian square
\[
\begin{CD}
 X'@>{g'}>>X\\
 @V{f'}VV @VV{f}V\\
 Y'@>{g}>>Y
\end{CD}
\tag{\(\square\)}
\]
with \(f\) separated and of finite type.  For a dense compactification
\(c=(j:X\hookrightarrow\overline X,\ p:\overline X\to Y)\), put
\[
 \overline X^\square=\overline X\times_YY'
\]
and denote the projections and induced maps by
\[
 \bar g:\overline X^\square\to\overline X,\qquad
 j^\square:X'\hookrightarrow\overline X^\square,\qquad
 p^\square:\overline X^\square\to Y'.
\]
The pair \(c^\square=(j^\square,p^\square)\) is a compactification of
\(f'\).  No density assertion is imposed on \(j^\square\).
\end{definition}

\begin{definition}[Dense closure after base change]
Let \(\overline X^{\square,\mathrm{dens}}\) be the scheme-theoretic closure
of \(X'\) in \(\overline X^\square\).  It defines a dense compactification
\(c^{\square,\mathrm{dens}}\) of \(f'\) and a proper refinement
\[
 \delta_c:c^{\square,\mathrm{dens}}\longrightarrow c^\square.
\]
The refinement \(\delta_c\) is Cartesian over \(X'\), and the associated
transformation \(\theta_{\delta_c}\) is defined.
\end{definition}

\begin{theorem}[Raw base change and dense replacement]
The morphism \(j^\square\) is a quasi-compact open immersion and
\(p^\square\) is proper.  The dense closure is proper over \(Y'\), and the
map
\[
 \theta_{\delta_c}:
 Rf'_{!,c^{\square,\mathrm{dens}}}
 \xrightarrow{\sim}Rf'_{!,c^\square}
\]
is a natural isomorphism.
\end{theorem}
\pred{stability of open and proper morphisms under base change, dense replacement, and Cartesian-refinement coherence}

\begin{definition}[The raw compact-support exchange map]
For \(K\in D^+(X,\Lambda_n)\), define
\[
 \operatorname{BC}^{c,\mathrm{raw}}_{!,\square}(K):
 g^*Rf_{!,c}K\longrightarrow Rf'_{!,c^\square}g'^*K
\]
as the composite
\[
\begin{aligned}
 g^*Rp_*j_!K
 &\xrightarrow{\operatorname{BC}_{p,g}(j_!K)}
 Rp^\square_*\bar g^*j_!K\\
 &\xrightarrow{Rp^\square_*\beta_{j,\bar g}(K)}
 Rp^\square_*j^\square_!g'^*K.
\end{aligned}
\]
Every source, target, and intermediate object lies in the indicated ambient
derived category.
\end{definition}

\begin{definition}[The raw inverse candidate]
The inverse candidate is
\[
\begin{aligned}
 \operatorname{IC}^{c,\mathrm{raw}}_{!,\square}(K):
 Rp^\square_*j^\square_!g'^*K
 &\xrightarrow{Rp^\square_*\beta_{j,\bar g}(K)^{-1}}
 Rp^\square_*\bar g^*j_!K\\
 &\xrightarrow{\operatorname{BC}_{p,g}(j_!K)^{-1}}
 g^*Rp_*j_!K.
\end{aligned}
\]
\end{definition}

\begin{lemma}[Left inverse equation]
For every \(K\in D^+(X,\Lambda_n)\),
\[
 \operatorname{IC}^{c,\mathrm{raw}}_{!,\square}(K)\,
 \operatorname{BC}^{c,\mathrm{raw}}_{!,\square}(K)
 =\operatorname{id}_{g^*Rf_{!,c}K}.
\]
\end{lemma}
\pred{the raw exchange map, its inverse candidate, and the inverse equations for proper and open exchange}

\begin{lemma}[Right inverse equation]
For every \(K\in D^+(X,\Lambda_n)\),
\[
 \operatorname{BC}^{c,\mathrm{raw}}_{!,\square}(K)\,
 \operatorname{IC}^{c,\mathrm{raw}}_{!,\square}(K)
 =\operatorname{id}_{Rf'_{!,c^\square}g'^*K}.
\]
\end{lemma}
\pred{the raw exchange map, its inverse candidate, and the inverse equations for proper and open exchange}

\begin{theorem}[Raw finite compact-support base change]
The transformation
\[
 \operatorname{BC}^{c,\mathrm{raw}}_{!,\square}:
 g^*Rf_{!,c}\xrightarrow{\sim}Rf'_{!,c^\square}g'^*
\]
is a natural isomorphism of functors
\(D^+(X,\Lambda_n)\to D^+(Y',\Lambda_n)\).
\end{theorem}
\pred{the left and right inverse equations}

\section{Refinement independence and canonical base change}

\begin{definition}[Base change of a refinement]
For a Cartesian refinement \(q:c_1\to c_2\), let
\[
 q^\square=q\times_YY':
 \overline X_{c_1}\times_YY'\longrightarrow
 \overline X_{c_2}\times_YY'
\]
be the induced refinement \(c^\square_1\to c^\square_2\).  Its source
square is the base change of the source square for \(q\), so it is
Cartesian.  Its refinement transformation is denoted
\(\theta_{q^\square}\).
\end{definition}

\begin{theorem}[Refinement naturality of compact-support exchange]
For every Cartesian refinement \(q:c_1\to c_2\), the square of natural
transformations
\[
\begin{CD}
 g^*Rf_{!,c_1}@>{\operatorname{BC}^{c_1,\mathrm{raw}}_{!,\square}}>>
 Rf'_{!,c^\square_1}g'^*\\
 @V{g^*\theta_q}VV @VV{\theta_{q^\square}g'^*}V\\
 g^*Rf_{!,c_2}@>{\operatorname{BC}^{c_2,\mathrm{raw}}_{!,\square}}>>
 Rf'_{!,c^\square_2}g'^*
\end{CD}
\]
commutes.
\end{theorem}
\pred{refinement transformations, naturality and pasting for proper base change, and naturality and pasting for open exchange}

\begin{corollary}[Independence under common refinement]
If \(d\) is a common refinement of \(c_1,c_2\), conjugating
\(\operatorname{BC}^{d,\mathrm{raw}}_{!,\square}\) by
\(\kappa_{c_1,d}\) and \(\kappa_{d,c_2}\) gives both
\(\operatorname{BC}^{c_1,\mathrm{raw}}_{!,\square}\) and
\(\operatorname{BC}^{c_2,\mathrm{raw}}_{!,\square}\), after inserting the
comparisons between the raw base-changed compactifications.  Hence the exchange map
depends only on \(f\) and the Cartesian square.
\end{corollary}
\pred{refinement naturality and independence of common-refinement comparisons}

\begin{definition}[Canonical compact-support exchange]
Let \(c_f\) and \(c_{f'}\) be the chosen dense compactifications defining
\(Rf_!\) and \(Rf'_!\).  Define
\[
 \operatorname{BC}_{!,\square}:
 g^*Rf_!\longrightarrow Rf'_!g'^*
\]
as
\[
 g^*Rf_{!,c_f}
 \xrightarrow{\operatorname{BC}^{c_f,\mathrm{raw}}_{!,\square}}
 Rf'_{!,c_f^\square}g'^*
 \xrightarrow{\kappa_{c_f^\square,c_{f'}}\,g'^*}
 Rf'_{!,c_{f'}}g'^*.
\]
The second arrow may equivalently be computed through the dense closure
\(c_f^{\square,\mathrm{dens}}\) and \(\theta_{\delta_{c_f}}\).
\end{definition}

\begin{theorem}[Finite-torsion compact-support base change]
The canonical transformation is a natural isomorphism
\[
 \operatorname{BC}_{!,\square}:
 g^*Rf_!\xrightarrow{\sim}Rf'_!g'^*
\]
of functors \(D^+(X,\Lambda_n)\to D^+(Y',\Lambda_n)\), independent of every
compactification and refinement used in its construction.
\end{theorem}
\pred{raw compact-support base change and independence under common refinement}

\section{Composition and units for the ambient functor}

\begin{definition}[Successive compactification datum]
Let \(X\xrightarrow{f}Y\xrightarrow{h}Z\), and choose compactifications
\[
 X\xrightarrow{j_f}\overline X_f\xrightarrow{p_f}Y,
 \qquad
 Y\xrightarrow{j_h}\overline Y_h\xrightarrow{p_h}Z.
\]
A successive compactification datum is a Cartesian square
\[
\begin{CD}
 \overline X_f@>{a}>>W\\
 @V{p_f}VV @VV{b}V\\
 Y@>{j_h}>>\overline Y_h
\end{CD}
\]
with \(a\) a quasi-compact schematically dense open immersion and \(b\)
proper.  It determines the compactification
\[
 d_{h,f}=\bigl(X\xrightarrow{aj_f}W\xrightarrow{p_hb}Z\bigr)
\]
of \(hf\).
\end{definition}

\begin{theorem}[Existence of successive compactifications]
Every pair of compactifications of \(f\) and \(h\) admits a successive
compactification datum.  Such data admit common dense refinements
compatible with the compactifications of \(f\), \(h\), and \(hf\).
\end{theorem}
\pred{Nagata compactification, scheme-theoretic closure, and common dense refinements}

\begin{definition}[Raw and canonical compositors]
For a successive compactification datum, define
\[
\begin{aligned}
 \widetilde c_{h,f}:Rh_{!,c_h}Rf_{!,c_f}
 &=Rp_{h,*}j_{h,!}Rp_{f,*}j_{f,!}\\
 &\xrightarrow{Rp_{h,*}\operatorname{Ex}_{j_h,b}}
 Rp_{h,*}Rb_*a_!j_{f,!}\\
 &\xrightarrow{\mu\text{ and open composition}}
 R(p_hb)_*(aj_f)_!=R(hf)_{!,d_{h,f}}.
\end{aligned}
\]
If \(c_{hf}\) is the chosen compactification defining \(R(hf)_!\), set
\[
 c_{h,f}:=
 \kappa_{d_{h,f},c_{hf}}\,\widetilde c_{h,f}:
 Rh_!Rf_!\xrightarrow{\sim}R(hf)_!.
\]
Both are natural isomorphisms of functors
\(D^+(X,\Lambda_n)\to D^+(Z,\Lambda_n)\).
\end{definition}

\begin{theorem}[Independence of the compositor]
The map \(c_{h,f}\) is independent of the compactifications and successive
compactification datum after insertion of the comparison maps \(\kappa\).
It is natural in the coefficient complex.
\end{theorem}
\pred{open--proper interchange, common dense refinements, refinement coherence, and proper-direct-image composition}

\begin{definition}[Identity unit]
The identity compactification
\(X\xrightarrow{\operatorname{id}}X\xrightarrow{\operatorname{id}}X\)
defines
\[
 u_X:R(\operatorname{id}_X)_!\xrightarrow{\sim}
 \operatorname{id}_{D^+(X,\Lambda_n)}.
\]
\end{definition}

\begin{theorem}[Unit and associativity equations]
For every compactifiable \(f:X\to Y\), the two composites formed from
\(c_{f,\operatorname{id}}\), \(c_{\operatorname{id},f}\), \(u_X\), and
\(u_Y\) are identities.  For three composable compactifiable morphisms
\(X\xrightarrow{f}Y\xrightarrow{h}Z\xrightarrow{k}W\),
\[
 c_{k,hf}\circ(Rk_!c_{h,f})
 =c_{kh,f}\circ(c_{k,h}Rf_!)
\]
as natural transformations \(Rk_!Rh_!Rf_!\to R(khf)_!\).
\end{theorem}
\pred{independence of the compositor, the cocycle for compactification comparisons, and associativity of open and proper composition}

\section{Pasting and pseudofunctorial coherence}

\begin{theorem}[Identity-square normalization]
If \(\square\) is the Cartesian square formed with identity maps, then
\(\operatorname{BC}_{!,\square}\) is the identity transformation after the
fixed unit identifications for inverse image and \(Rf_!\).
\end{theorem}
\pred{the identity equations for open exchange, proper base change, and refinement comparisons}

\begin{theorem}[Vertical pasting]
For two vertically composable Cartesian squares, the base-change
isomorphism for the outer rectangle is the composite of the base-change
isomorphisms for the two squares, with the canonical pullback associator
inserted.  Explicitly, if base change by \(g_1\) is followed by base change
by \(g_2\), then
\[
 \operatorname{BC}_{!,g_1g_2}
 =
 \bigl(\operatorname{BC}_{!,g_2}\,g_1'^*\bigr)
 \circ
 \bigl(g_2^*\operatorname{BC}_{!,g_1}\bigr).
\]
\end{theorem}
\pred{vertical pasting for open exchange and proper base change, and refinement naturality}

\begin{theorem}[Horizontal pasting]
For two horizontally composable Cartesian squares whose vertical maps are
compactifiable, the base-change isomorphism for the outer rectangle agrees
with the composite formed from the two base-change maps and the compositor
\[
 c_{h,f}:Rh_!Rf_!\xrightarrow{\sim}R(hf)_!.
\]
The comparison is unchanged by replacing any compactification used to
define either compositor or base-change map.
\end{theorem}
\pred{horizontal pasting for open exchange and proper base change, refinement naturality, and compactification-independent compositors}

\begin{corollary}[Beck--Chevalley coherence]
The functors \(g^*\) and \(Rf_!\), their units and compositors, and the maps
\(\operatorname{BC}_{!,\square}\) form a coherent Beck--Chevalley system on
the category of the stated schemes and compactifiable morphisms.  All
diagrams obtained from identity squares and finite rectangular pastings
commute.
\end{corollary}
\pred{identity normalization and vertical and horizontal pasting}

\section{Fibers and compactly supported cohomology}

\begin{definition}[Compactly supported cohomology]
Let \(k\) be a field, let \(a:X\to\operatorname{Spec}k\) be separated and
of finite type, and let \(K\in D^+(X,\Lambda_n)\).  Define
\[
 R\Gamma_c(X,K):=Ra_!K,\qquad
 H_c^q(X,K):=\mathcal H^q(Ra_!K).
\]
For a geometric point \(\bar y\to Y\), write
\(X_{\bar y}=X\times_Y\bar y\) and let
\(a_{\bar y}:X_{\bar y}\to\bar y\) be its structure morphism.
\end{definition}

\begin{corollary}[Compact-support stalk formula]
For \(f:X\to Y\), \(K\in D^+(X,\Lambda_n)\), and a geometric point
\(\bar y\to Y\), compact-support base change gives
\[
 (Rf_!K)_{\bar y}\simeq
 R\Gamma_c(X_{\bar y},K|_{X_{\bar y}}),
\]
and hence
\[
 (R^qf_!K)_{\bar y}\simeq
 H_c^q(X_{\bar y},K|_{X_{\bar y}}).
\]
\end{corollary}
\pred{finite-torsion compact-support base change and the definition of compactly supported cohomology}

\begin{corollary}[Ground-field extension]
For a field extension \(v:\operatorname{Spec}k'\to\operatorname{Spec}k\),
let \(X'=X\times_kk'\), with projection \(v':X'\to X\) and structure
morphism \(a':X'\to\operatorname{Spec}k'\).  Then
\[
 v^*Ra_!K\xrightarrow{\sim}Ra'_!v'^*K
\]
in \(D(\operatorname{Spec}k',\Lambda_n)\).  If
\(\Lambda_n\to\Lambda'_n\) is a finite coefficient-ring extension of
finite Tor dimension, then
\[
 v^*Ra_!K\otimes_{\Lambda_n}^{L}\Lambda_n'
 \xrightarrow{\sim}
 Ra'_!(v'^*K\otimes_{\Lambda_n}^{L}\Lambda_n').
\]
\end{corollary}
\pred{finite-torsion compact-support base change and scalar-extension compatibility}

\section{The finite-level finiteness chain}

\begin{definition}[Quasi-finite seed]
The quasi-finite seed is the assertion that a separated quasi-finite
finite-type morphism \(u:V\to T\) has cohomological amplitude \([0,0]\) for
exceptional direct image on finite constructible sheaves and that
\(u_!\) preserves finite constructibility.
\end{definition}

\begin{theorem}[Quasi-finite exceptional image]
The quasi-finite seed holds.  More precisely, a Zariski-main factorization
\[
 V\xrightarrow{j}\overline V\xrightarrow{\pi}T
\]
with \(j\) open and \(\pi\) finite identifies
\(Ru_!\) with the exact functor \(\pi_*j_!\).  Thus
\[
 Ru_!:\Dbc(V,\Lambda_n)\longrightarrow\Dbc(T,\Lambda_n).
\]
\end{theorem}
\pred{Zariski's main theorem, exact open extension, and exact finite direct image}

\begin{definition}[Relative affine-line test]
For a noetherian \(T\), an open immersion
\(j:U\hookrightarrow\mathbf A^1_T\), the projection
\(\pi:U\to T\), and a constructible sheaf \(F\) on \(U\), let
\(\operatorname{AL}(T,U,F)\) denote the assertion that every
\(R^q\pi_!F\) is constructible and is zero for \(q\notin[0,2]\), compatibly
with arbitrary base change on \(T\).
\end{definition}

\begin{theorem}[Relative-curve constant-coefficient seed]
The assertion \(\operatorname{AL}(T,U,\Lambda_n)\) holds for the constant
finite coefficient sheaf.  The same statement holds for every finite
constant \(\Lambda_n\)-module and is stable under finite direct sums,
kernels, cokernels, and extensions.
\end{theorem}
\pred{compact-support stalks on affine curves, finite fiber cohomology, and the open-extension exact sequence on \(\mathbf P^1_T\)}

\begin{definition}[Finite lisse descent datum]
For a finite locally constant \(\Lambda_n\)-sheaf \(L\) on \(U\), a finite
lisse descent datum consists of a finite etale surjection
\(h:V\to U\), a finite \(\Lambda_n\)-module \(M\), an isomorphism
\[
 h^*L\xrightarrow{\sim}M_V,
\]
and the induced descent isomorphism between the two pullbacks of \(M_V\) to
\(V\times_UV\), satisfying the identity and triple-overlap cocycle
equations.
\end{definition}

\begin{theorem}[Finite monodromy and finite etale trivialization]
Every finite locally constant \(\Lambda_n\)-sheaf admits a finite lisse
descent datum.  Morphisms of such sheaves are equivalent to morphisms of
their descent data after passage to one common finite etale refinement.
\end{theorem}
\pred{the finite automorphism group of a finite \(\Lambda_n\)-module and effective descent for finite etale sheaves}

\begin{theorem}[Relative-curve finiteness for lisse coefficients]
For every finite locally constant \(\Lambda_n\)-sheaf \(L\) on an open
subscheme \(U\subset\mathbf A^1_T\), the assertion
\(\operatorname{AL}(T,U,L)\) holds.  The isomorphisms supplied by arbitrary
base change on \(T\) commute with the descent cocycle of every finite lisse
descent datum.
\end{theorem}
\pred{finite etale trivialization, finite etale cohomological descent for compact support, and the constant-coefficient seed}

\begin{theorem}[Restriction to an open part of the base]
Let \(v:T_0\hookrightarrow T\) be open, let
\(U_0=U\times_TT_0\), and let \(\pi_0:U_0\to T_0\).  For every
constructible \(F\) on \(U\), canonical compact-support base change gives
\[
 v^*R\pi_!F\xrightarrow{\sim}
 R\pi_{0,!}(F|_{U_0}).
\]
Thus constructibility and a fixed cohomological bound on \(T_0\) may be
tested after restriction to this open part.
\end{theorem}
\pred{finite-torsion compact-support base change}

\begin{theorem}[Closed-base localization and constructible gluing]
Let \(v:T_0\hookrightarrow T\) be open and let
\(w:T_1\hookrightarrow T\) be its complementary reduced closed immersion.
For \(M\in D(T,\Lambda_n)\), there is a functorial distinguished triangle
\[
 v_!v^*M\longrightarrow M\longrightarrow
 Rw_*w^*M\longrightarrow.
\]
The object \(M\) lies in \(\Dbc(T,\Lambda_n)\) if and only if
\(v^*M\in\Dbc(T_0,\Lambda_n)\) and
\(w^*M\in\Dbc(T_1,\Lambda_n)\), with finite bounds on both pieces.
For \(M=R\pi_!F\), compact-support base change identifies these restrictions
with
\[
 R\pi_{0,!}(F|_{U_0}),\qquad
 R\pi_{1,!}(F|_{U_1}),
\]
where \(U_i=U\times_TT_i\).
\end{theorem}
\pred{the open--closed localization triangle, proper direct image for a closed immersion, and finite-torsion compact-support base change}

\begin{definition}[Affine-line devissage datum]
For \((T,U,F)\), an affine-line devissage datum consists of a dense open
\(v:T_0\hookrightarrow T\), an open immersion
\(a:W\hookrightarrow U_0=U\times_TT_0\), and its closed complement
\(i:Z\hookrightarrow U_0\), such that \(F|_W\) is finite locally constant
and \(Z\to T_0\) is quasi-finite.  The remaining closed base is
\(T_1=T\setminus T_0\), equipped with its reduced induced structure.
\end{definition}

\begin{theorem}[Existence of affine-line devissage data]
For every constructible finite \(\Lambda_n\)-sheaf on
\(U\subset\mathbf A^1_T\), there is a finite chain of reduced closed
subschemes
\[
 T=T^{(0)}\supsetneq T^{(1)}\supsetneq\cdots
 \supsetneq T^{(N)}=\varnothing
\]
such that the restriction over each locally closed layer
\(T^{(r)}\setminus T^{(r+1)}\) admits an affine-line devissage datum.
\end{theorem}
\pred{constructible stratifications, generic finiteness of a proper closed subset of a relative affine line, and noetherianity of \(T\)}

\begin{proposition}[One devissage step]
Let \((T_0,W,Z,T_1)\) be an affine-line devissage datum.  Assume
\(\operatorname{AL}(T_0,W,F|_W)\), the quasi-finite exceptional-image
theorem for \(Z\to T_0\), and the assertion
\(\operatorname{AL}(T_1,U\times_TT_1,F|_{T_1})\).  Then
\(\operatorname{AL}(T,U,F)\) holds.  The comparison is governed by the
distinguished triangle
\[
 a_!a^*(F|_{U_0})\longrightarrow F|_{U_0}
 \longrightarrow i_*i^*(F|_{U_0})\longrightarrow
\]
on \(U_0\), together with the restriction isomorphism over \(T_0\).
The closed-base localization triangle then glues the result over \(T_0\)
to the assumed result over \(T_1\).
\end{proposition}
\pred{relative-curve finiteness for lisse coefficients, quasi-finite exceptional image, restriction to an open base, closed-base localization, and the open--closed triangle}

\begin{theorem}[Open--closed devissage on the affine line]
For every constructible finite \(\Lambda_n\)-sheaf \(F\) on
\(U\subset\mathbf A^1_T\), the assertion
\(\operatorname{AL}(T,U,F)\) holds.
\end{theorem}
\pred{existence of affine-line devissage data, one devissage step, and noetherian induction on the closed base}

\begin{proposition}[Affine-space reduction]
For the projection \(\mathbf A^d_T\to T\), factorization into \(d\)
affine-line projections and the compositor for \(R(-)_!\) give
constructibility preservation and cohomological amplitude contained in
\([0,2d]\).
\end{proposition}
\pred{open--closed devissage on the affine line and composition of exceptional direct images}

\begin{proposition}[Local factorization of a separated finite-type morphism]
Zariski locally on \(Y\), a separated finite-type
\(f:X\to Y\) admits a finite affine cover on which the restricted morphisms
factor as
\[
 X_a\xrightarrow{i_a}\mathbf A^{d_a}_Y
 \xrightarrow{\pi_a}Y,
\]
with \(i_a\) a locally closed immersion.  The exceptional direct image for
\(i_a\) is controlled by the quasi-finite seed, and the exceptional direct
image for \(\pi_a\) is controlled by affine-space reduction.  The
Mayer--Vietoris triangles for the finite cover reconstruct \(Rf_!\).
\end{proposition}
\pred{separated finite-type affine embedding, the quasi-finite theorem, affine-space reduction, and finite Mayer--Vietoris triangles}

\begin{theorem}[Finite cohomological dimension]
If all geometric fibers of \(f:X\to Y\) have dimension at most \(d\), then
for every constructible sheaf \(F\),
\[
 R^qf_!F=0\qquad(q>2d).
\]
The bound is local on \(Y\), stable under base change, and subadditive under
composition.
\end{theorem}
\pred{the compact-support stalk formula, affine-space reduction, and the geometric cohomological-dimension bound}

\begin{theorem}[Constructibility and boundedness of
\texorpdfstring{\(Rf_!\)}{Rf!}]
For every separated finite-type \(f:X\to Y\),
\[
 Rf_!:\Dbc(X,\Lambda_n)\longrightarrow\Dbc(Y,\Lambda_n)
\]
is an exact triangulated functor.  If \(K\) has cohomology in \([a,b]\) and
the fibers of \(f\) have dimension at most \(d\), then \(Rf_!K\) has
cohomology in \([a,b+2d]\).
\end{theorem}
\pred{local factorization, finite cohomological dimension, and finite truncation triangles}

\begin{corollary}[Proper direct image finiteness]
For proper \(f\), the identity compactification identifies
\(Rf_!\simeq Rf_*\).  Hence
\[
 Rf_*:\Dbc(X,\Lambda_n)\longrightarrow\Dbc(Y,\Lambda_n).
\]
\end{corollary}
\pred{constructibility and boundedness of \(Rf_!\) and the proper normalization}

\begin{corollary}[Internal form of compact-support base change]
After the finiteness theorem, the canonical isomorphism
\[
 g^*Rf_!\xrightarrow{\sim}Rf'_!g'^*
\]
is an isomorphism of functors
\(\Dbc(X,\Lambda_n)\to\Dbc(Y',\Lambda_n)\).
\end{corollary}
\pred{finite-torsion compact-support base change and constructibility and boundedness of \(Rf_!\)}

\section{Integral adic transport}

\begin{convention}[Continuous integral coefficients on the pro-etale site]
Let \(E/\Ql\) be finite, let \(\mathcal O_E\) be its ring of integers, let
\(\varpi\) be a uniformizer, and put
\(\Lambda_m=\mathcal O_E/\varpi^m\).  From this section onward the site is
\(T_{\mathrm{pro\acute et}}\).  If \((\Lambda_m)_T\) denotes the finite
constant sheaf of rings on this site, define the continuous coefficient
sheaf
\[
 (\mathcal O_E)_T:=\varprojlim_m(\Lambda_m)_T
\]
in sheaves of rings.  Write \(\mathcal O_E\) for this ring sheaf.  Let
\(\Mod_{\mathrm{pro\acute et}}(T,\mathcal O_E)\) be the Grothendieck abelian category of
sheaves of modules over \((\mathcal O_E)_T\), and put
\[
 D(T,\mathcal O_E)=D\bigl(\Mod_{\mathrm{pro\acute et}}(T,\mathcal O_E)\bigr).
\]
All restriction-of-scalars maps, tensor products, and completeness
conditions below are taken over this ring object.  Derived tensor products
are computed by K-flat replacements in its module-sheaf category.
\end{convention}

\begin{definition}[Strict coefficient systems]
The abelian category
\(\Mod_{\mathrm{pro\acute et}}(T,\mathcal O_{E,\bullet})\) has systems
\[
 M_\bullet=(M_m,\rho_{m+1,m})_{m\ge1},
\]
where \(M_m\) is a \(\Lambda_m\)-module sheaf and the transition maps are
\(\mathcal O_E\)-linear after restriction of scalars.  Kernels and
cokernels are levelwise.  Let
\[
 D(T,\mathcal O_{E,\bullet})=
 D\bigl(\Mod_{\mathrm{pro\acute et}}(T,\mathcal O_{E,\bullet})\bigr),
\]
with exact evaluation functors \(\operatorname{ev}_m\).
\end{definition}

\begin{definition}[Normalized bounded constructible system]
A system \(K_\bullet\in D(T,\mathcal O_{E,\bullet})\) is normalized if
\[
 \operatorname{ev}_{m+1}(K_\bullet)
 \otimes^L_{\Lambda_{m+1}}\Lambda_m
 \xrightarrow{\sim}\operatorname{ev}_m(K_\bullet)
\]
coherently for all \(m\).  It has uniform bounded constructible cohomology
if every level lies in \(\Dbc(T,\Lambda_m)\) and one finite interval bounds
all levels.
\end{definition}

\begin{theorem}[Normalized cone theorem]
Normalized systems with uniform bounded constructible cohomology are
closed under shifts.  The cone of any morphism between two such systems is
normalized, its level \(m\) is the cone of the level-\(m\) morphism, and it
has one uniform cohomological bound.  Hence they form a full triangulated
subcategory
\[
 \Dbcnorm(T,\mathcal O_{E,\bullet})
 \subset D(T,\mathcal O_{E,\bullet}).
\]
\end{theorem}
\pred{exact evaluation, derived tensor compatibility with cones, and the two-out-of-three property}

\begin{definition}[Towers, derived inverse limit, and realization]
Define one abelian tower category
\[
 \mathsf{Tow}_T(\mathcal O_E)=
 \operatorname{Fun}(\mathbf N_{\ge1}^{\mathrm{op}},
                    \Mod_{\mathrm{pro\acute et}}(T,\mathcal O_E)).
\]
Levelwise restriction of scalars gives an exact functor
\[
 \operatorname{Res}_T:
 D(T,\mathcal O_{E,\bullet})
 \longrightarrow D(\mathsf{Tow}_T(\mathcal O_E)).
\]
The left exact inverse limit has a right-derived functor, constructed by a
K-injective resolution in this single tower category,
\[
 R\!\varprojlim:
 D(\mathsf{Tow}_T(\mathcal O_E))\to D(T,\mathcal O_E).
\]
Set
\(\operatorname{Real}_T=R\!\varprojlim\circ\operatorname{Res}_T\).
The tower projections define
\[
 r_m(K_\bullet):
 \operatorname{Real}_T(K_\bullet)
 \otimes^L_{\mathcal O_E}\Lambda_m
 \longrightarrow\operatorname{ev}_m(K_\bullet).
\]
\end{definition}

\begin{theorem}[Adic reconstruction and conservativity]
For a normalized bounded constructible system, every \(r_m\) is an
isomorphism and the realization is derived \(\varpi\)-complete.  A morphism
between two such realizations is an isomorphism if and only if all its
finite reductions are isomorphisms.
\end{theorem}
\pred{the quotient-tower Mittag--Leffler property, derived completion, and the maps \(r_m\)}

\begin{definition}[Integral constructible category]
The category \(\Dbc(T,\mathcal O_E)\) is the full subcategory of
\(D(T,\mathcal O_E)\) of derived \(\varpi\)-complete objects \(K\) for which
\[
 \operatorname{red}_mK=
 K\otimes^L_{\mathcal O_E}\Lambda_m
\]
forms a normalized bounded constructible system with one uniform bound.
Adic reconstruction gives its ordinary triangulated structure.
\end{definition}

\begin{definition}[Completed inverse image]
For \(h:T'\to T\) and \(K\in\Dbc(T,\mathcal O_E)\), define
\[
 L\widehat h^{\,*}K:=
 \operatorname{Real}_{T'}
 \bigl((h^*\operatorname{red}_mK)_{m\ge1}\bigr).
\]
This is the completed inverse image; ordinary termwise pullback is followed
by derived-complete realization in one \(\mathcal O_E\)-module category.
\end{definition}

\begin{theorem}[Finite reductions of completed inverse image]
Completed inverse image is exact, preserves bounded constructibility, and
satisfies
\[
 \operatorname{red}_m(L\widehat h^{\,*}K)
 \xrightarrow{\sim}h^*(\operatorname{red}_mK).
\]
It has coherent identity and composition isomorphisms.
\end{theorem}
\pred{finite-level pullback, normalized cones, and adic reconstruction}

\begin{theorem}[Levelwise exceptional direct image]
Finite-level exceptional direct image defines an exact functor
\[
 Rf_{!,\bullet}:
 \Dbcnorm(X,\mathcal O_{E,\bullet})
 \longrightarrow\Dbcnorm(Y,\mathcal O_{E,\bullet})
\]
whose level \(m\) is \(Rf_!\).  Realization gives
\[
 Rf_!:\Dbc(X,\mathcal O_E)\longrightarrow\Dbc(Y,\mathcal O_E),
\qquad
 \operatorname{red}_mRf_!K\simeq Rf_!\operatorname{red}_mK.
\]
\end{theorem}
\pred{finite-level constructibility and boundedness, reduction compatibility of open and proper direct image, and adic reconstruction}

\begin{theorem}[Reduction compatibility of compact-support exchange]
For \(K\in\Dbc(X,\mathcal O_E)\), the level-\(m\) reduction of the adic
exchange candidate is the finite-torsion isomorphism
\[
 \operatorname{BC}_{!,\square}(\operatorname{red}_mK):
 g^*Rf_!(\operatorname{red}_mK)
 \xrightarrow{\sim}
 Rf'_!g'^*(\operatorname{red}_mK).
\]
These maps commute with all transitions \(\Lambda_{m+1}\to\Lambda_m\).
\end{theorem}
\pred{finite-level internal compact-support base change, completed pullback reductions, and levelwise exceptional direct image}

\begin{definition}[Integral adic compact-support exchange]
The compatible finite-level maps define, after realization,
\[
 \operatorname{BC}^{\mathcal O_E}_{!,\square}:
 L\widehat g^{\,*}Rf_!
 \longrightarrow
 Rf'_!L\widehat {g'}^{\,*}.
\]
\end{definition}

\begin{theorem}[Integral adic compact-support base change]
The integral adic exchange is a natural isomorphism
\[
 L\widehat g^{\,*}Rf_!
 \xrightarrow{\sim}
 Rf'_!L\widehat {g'}^{\,*}
\]
on \(\Dbc(X,\mathcal O_E)\).  Its reduction at every level is the
finite-torsion exchange.  Identity, refinement, composition, and pasting
equations hold and are detected by finite reductions.
\end{theorem}
\pred{reduction compatibility, adic reconstruction, and conservativity of all finite reductions}

\section{Rational and algebraically closed transport}

\begin{definition}[Central localization]
Let
\[
 \mathcal A_E(T)=\Dbc(T,\mathcal O_E)[\varpi^{-1}]
\]
have the same objects as the integral category and morphisms
\[
 \Hom_{\mathcal A_E(T)}(K,L)=
 \Hom_{\Dbc(T,\mathcal O_E)}(K,L)\otimes_{\mathcal O_E}E.
\]
Its shift is inherited from the integral category, and a triangle is
distinguished if it is isomorphic to the localized image of an integral
distinguished triangle.
\end{definition}

\begin{definition}[Karoubi envelope]
The Karoubi envelope \(\Kar(\mathcal A_E(T))\) has objects \((K,e)\),
where \(K\in\mathcal A_E(T)\) and
\(e^2=e\in\operatorname{End}_{\mathcal A_E(T)}(K)\).  A morphism
\((K,e)\to(L,d)\) is a morphism \(u:K\to L\) satisfying \(u=due\), and
the shift is \((K,e)[1]=(K[1],e[1])\).
\end{definition}

\begin{theorem}[Central localization and Karoubi triangulation]
The category \(\mathcal A_E(T)\) is \(E\)-linear triangulated.  Its
Karoubi envelope has the triangulated idempotent-completion structure in
which every distinguished triangle is a direct summand of the image of a
distinguished triangle in \(\mathcal A_E(T)\).  Define
\[
 \Dbc(T,E)=\Kar(\mathcal A_E(T)).
\]
\end{theorem}
\pred{triangulated localization by the central action of \(\varpi\) and the triangulated idempotent-completion theorem}

\begin{theorem}[Idempotent transport]
Every exact integral functor extends through central localization and then
acts on idempotent objects by
\[
 (K,e)\longmapsto(FK,F(e)).
\]
A natural isomorphism \(\alpha:F\to G\) restricts to
\[
 G(e)\,\alpha_K\,F(e):(FK,F(e))\longrightarrow(GK,G(e)).
\]
This transport preserves distinguished triangles, identities,
composition, and equations between natural transformations.
\end{theorem}
\pred{central localization and Karoubi triangulation}

\begin{definition}[Finite scalar extension]
For a finite extension \(E'/E\), the continuous ring map
\((\mathcal O_E)_T\to(\mathcal O_{E'})_T\) defines
\[
 K\otimes_EE':=
 K\otimes^L_{(\mathcal O_E)_T}(\mathcal O_{E'})_T.
\]
After central localization and idempotent completion, set
\[
 (K,e)\otimes_EE'=
 (K\otimes_EE',e\otimes 1).
\]
This formula defines a functor
\(\Dbc(T,E)\to\Dbc(T,E')\).
\end{definition}

\begin{theorem}[Exactness and coherence of finite scalar extension]
Finite scalar extension preserves derived completeness, bounded
constructibility, distinguished triangles, and idempotent summands.  For
\(E\subset E'\subset E''\), its canonical associator identifies the two
functors from \(\Dbc(T,E)\) to \(\Dbc(T,E'')\), and these associators obey
the identity and triple-extension equations.
\end{theorem}
\pred{finite flatness of rings of integers, finite-reduction compatibility, and idempotent transport}

\begin{definition}[Rational exceptional direct image]
For a separated finite-type \(f:X\to Y\), idempotent transport of the
integral adic functor defines
\[
 Rf_{!,E}:\Dbc(X,E)\longrightarrow\Dbc(Y,E),
 \qquad
 (K,e)\longmapsto(Rf_!K,Rf_!(e)).
\]
We write this exact \(E\)-linear functor as \(Rf_!\).
\end{definition}

\begin{theorem}[Scalar compatibility of rational exceptional image]
For every finite extension \(E'/E\), there is a natural isomorphism
\[
 (Rf_{!,E}K)\otimes_EE'
 \xrightarrow{\sim}
 Rf_{!,E'}(K\otimes_EE').
\]
It is compatible with towers of finite extensions, compactification
comparisons, identity units, and compositors.
\end{theorem}
\pred{finite-reduction compatibility of integral exceptional direct image and idempotent transport}

\begin{definition}[Rational inverse image]
For \(h:T'\to T\), the exact integral functor
\(L\widehat h^{\,*}\) extends by idempotent transport to an exact
\(E\)-linear functor
\[
 h_E^*:\Dbc(T,E)\longrightarrow\Dbc(T',E).
\]
We write this functor as \(h^*\).  Its identity and composition
isomorphisms are those transported from completed inverse image.
\end{definition}

\begin{theorem}[Scalar compatibility of rational inverse image]
For every finite extension \(M/E\), there is a coherent natural
isomorphism
\[
 (h_E^*K)\otimes_EM\xrightarrow{\sim}
 h_M^*(K\otimes_EM).
\]
It is compatible with towers of finite extensions and with the identity
and composition isomorphisms for inverse image.
\end{theorem}
\pred{finite reductions of completed inverse image, exact finite scalar extension, and idempotent transport}

\begin{theorem}[Rational compact-support base change]
Localization and idempotent transport give a natural isomorphism
\[
 \operatorname{BC}^{E}_{!,\square}:
 g^*Rf_!\xrightarrow{\sim}Rf'_!g'^*
\]
on \(\Dbc(X,E)\).  It is independent of an integral lattice and commutes
with scalar extension along every finite extension \(E'/E\).
\end{theorem}
\pred{integral adic compact-support base change and idempotent transport}

\begin{definition}[Finite coefficient-field pair category]
Fix \(\Qlb\).  Let \(\mathcal C_c(T,\Qlb)\) have objects
\((E,K_E)\), with \(E\subset\Qlb\) finite over \(\Ql\) and
\(K_E\in\Dbc(T,E)\), and morphisms
\[
 \Hom((E,K_E),(E',L_{E'}))=
 \varinjlim_{M\supset E,E'}
 \Hom_{\Dbc(T,M)}
 (K_E\otimes_EM,L_{E'}\otimes_{E'}M).
\]
The transition maps in this filtered colimit are exact scalar-extension
functors.  Identities and composition are represented after extension to
one common finite coefficient field and are independent of a further finite extension.
\end{definition}

\begin{definition}[Common-stage shifts, cones, and triangles]
Set \((E,K_E)[1]=(E,K_E[1])\).  If a morphism is represented by
\(u_M:K_E\otimes_EM\to L_{E'}\otimes_{E'}M\), its common-stage cone is
\((M,\operatorname{Cone}(u_M))\).  A triangle in
\(\mathcal C_c(T,\Qlb)\) is distinguished if, after extension to one
finite coefficient field containing all fields of definition and all representatives,
it is represented by a distinguished triangle in \(\Dbc(T,M)\).
\end{definition}

\begin{theorem}[Independence of common-stage cones]
Two representatives of one morphism acquire isomorphic cones after a
further finite scalar extension.  The isomorphism class is compatible with
shifts, composition, and scalar extension.  Distinguishedness is
independent of the chosen common finite coefficient field.
\end{theorem}
\pred{filteredness of finite subextensions, exact scalar extension, and compatibility of mapping cones with exact functors}

\begin{theorem}[Triangulation of the pair category]
The pair category is a \(\Qlb\)-linear triangulated category, scalar
extension between finite coefficient fields is exact, and distinguishedness is
independent of the common finite coefficient field.
\end{theorem}
\pred{independence of common-stage cones and the triangulated axioms at one finite stage}

\begin{definition}[Algebraically closed coefficient category]
The algebraically closed constructible category is the triangulated
Karoubi envelope
\[
 \Dbc(T,\Qlb)=\Kar(\mathcal C_c(T,\Qlb)).
\]
An object is a pair \((A,e)\), with
\(A\in\mathcal C_c(T,\Qlb)\) and \(e^2=e\); morphisms, shifts, and
distinguished triangles are those of the triangulated idempotent
completion.
\end{definition}

\begin{definition}[Transport of a coherent finite coefficient-field family]
Let \(F_E:\Dbc(T,E)\to\Dbc(T',E)\) be exact functors, equipped for every
finite \(M/E\) with coherent isomorphisms
\[
 \sigma^F_{M/E,K}:(F_EK)\otimes_EM
 \xrightarrow{\sim}F_M(K\otimes_EM).
\]
Define \(F_{\mathrm{pair}}:\mathcal C_c(T,\Qlb)\to
\mathcal C_c(T',\Qlb)\) on objects by
\[
 F_{\mathrm{pair}}(E,K_E)=(E,F_EK_E).
\]
If \([M,u_M]\) represents a morphism
\((E,K_E)\to(E',L_{E'})\), define its image to be represented over \(M\)
by
\[
 (\sigma^F_{M/E',L_{E'}})^{-1}\,F_M(u_M)\,
 \sigma^F_{M/E,K_E}.
\]
\end{definition}

\begin{theorem}[Well-definedness of pair-level transport]
The preceding formula is independent of the common field and the
representative, preserves identities and composition, and defines an exact
triangulated functor.  Natural transformations of coherent finite coefficient-field
families induce natural transformations of pair-level functors by the
same common-stage formula.
\end{theorem}
\pred{scalar-extension coherence, independence of common-stage cones, and naturality after one finite extension}

\begin{definition}[Pair-level exceptional image and inverse image]
Apply pair-level transport to the coherent families
\((Rf_{!,E})_E\) and \((h_E^*)_E\).  Denote the resulting exact functors by
\[
 Rf_{!,\mathrm{pair}}:
 \mathcal C_c(X,\Qlb)\to\mathcal C_c(Y,\Qlb),
 \qquad
 h^*_{\mathrm{pair}}:
 \mathcal C_c(T,\Qlb)\to\mathcal C_c(T',\Qlb).
\]
\end{definition}

\begin{definition}[Pair-level compact-support exchange]
The scalar-compatible family \((\operatorname{BC}^E_{!,\square})_E\)
defines a natural transformation
\[
 \operatorname{BC}^{\mathrm{pair}}_{!,\square}:
 g^*_{\mathrm{pair}}Rf_{!,\mathrm{pair}}
 \longrightarrow
 Rf'_{!,\mathrm{pair}}g'^*_{\mathrm{pair}}.
\]
Its component on \((E,K_E)\) is represented at stage \(E\) by
\(\operatorname{BC}^E_{!,\square}(K_E)\); after passage to \(M/E\), it is
the conjugate of \(\operatorname{BC}^M_{!,\square}(K_E\otimes_EM)\) by the
four scalar-compatibility isomorphisms.
\end{definition}

\begin{theorem}[Pair-level compact-support base change]
The pair-level exchange is a natural isomorphism of exact triangulated
functors.  Its identity, refinement, compositor, and pasting equations may
be checked after passage to one common finite coefficient field.
\end{theorem}
\pred{rational compact-support base change, scalar compatibility, and well-defined pair-level transport}

\begin{definition}[Algebraically closed exceptional and inverse images]
For an object \((A,e)\) of the Karoubi envelope, define
\[
 Rf_!(A,e)=
 (Rf_{!,\mathrm{pair}}A,Rf_{!,\mathrm{pair}}(e)),
 \qquad
 h^*(A,e)=
 (h^*_{\mathrm{pair}}A,h^*_{\mathrm{pair}}(e)).
\]
These are exact functors on \(\Dbc(-,\Qlb)\).
\end{definition}

\begin{definition}[\texorpdfstring{\(\Qlb\)}{Qbar-l}-compact-support exchange]
For \((A,e)\in\Dbc(X,\Qlb)\), define
\[
 \operatorname{BC}^{\Qlb}_{!,\square}(A,e)=
 (Rf'_{!,\mathrm{pair}}g'^*_{\mathrm{pair}})(e)\,
 \operatorname{BC}^{\mathrm{pair}}_{!,\square}(A)\,
 (g^*_{\mathrm{pair}}Rf_{!,\mathrm{pair}})(e).
\]
This is the idempotent restriction of the pair-level transformation.
\end{definition}

\begin{theorem}[\texorpdfstring{\(\Qlb\)}{Qbar-l}-compact-support base change]
The idempotent restrictions define a natural isomorphism
\[
 \operatorname{BC}^{\Qlb}_{!,\square}:
 g^*Rf_!\xrightarrow{\sim}Rf'_!g'^*
\]
on \(\Dbc(X,\Qlb)\).  It is independent of the finite coefficient field of definition,
and its identity, refinement, compositor, and pasting equations hold after
passage to one common finite extension.
\end{theorem}
\pred{pair-level compact-support base change and transport to the triangulated Karoubi envelope}

\section{Bibliographic remarks}

The geometric ingredients are Nagata compactification, schematic closure,
proper base change, Zariski's main theorem, and compact-support finiteness
for relative curves.  The coefficient constructions use derived
completion in an abelian tower category, triangulated central localization,
and triangulated idempotent completion.  These results are developed in the
Stacks Project chapters on morphisms, etale and pro-etale cohomology,
derived completion, and triangulated categories.

\end{document}
